\kchapter{ReplicaSets and Deployments}{deploy}

\begin{objectives}
\begin{itemize}
  \item how a ReplicaSet keeps N copies of a Pod alive;
  \item why you use Deployments instead of ReplicaSets directly;
  \item rolling updates, \kc{maxSurge} and \kc{maxUnavailable};
  \item checking, pausing, and rolling back rollouts;
  \item release strategies: recreate, rolling, blue/green, canary.
\end{itemize}
\end{objectives}

\section{The ownership chain}

\begin{center}
\begin{tikzpicture}[
  obj/.style={draw=k8sblue, fill=white, line width=1pt, rounded corners=5pt,
              minimum width=3.0cm, minimum height=1.6cm, align=center,
              font=\small\bfseries, text=k8snavy},
  >={Stealth[length=3mm]}]
  \node[obj] (d) at (0,0) {\kimg[0.8cm]{deploy}\\Deployment\\[-2pt]\scriptsize\mdseries web};
  \node[obj] (r) at (4.4,0) {\kimg[0.8cm]{rs}\\ReplicaSet\\[-2pt]\scriptsize\mdseries web-7d9f8c};
  \foreach \y/\n in {1.5/xk2lp, 0/p9q4m, -1.5/b7c1z}
    \node[obj, minimum width=2.6cm, minimum height=1.0cm] (p\n) at (8.6,\y)
      {\kicon{pod}\ Pod\\[-2pt]\scriptsize\mdseries web-7d9f8c-\n};
  \draw[->, thick, k8sblue] (d) -- node[above, font=\scriptsize]{owns} (r);
  \foreach \n in {xk2lp, p9q4m, b7c1z} \draw[->, thick, k8sblue] (r) -- (p\n.west);
\end{tikzpicture}
\end{center}

\begin{itemize}
  \item A \kterm{ReplicaSet} has one job: ensure exactly \kc{replicas} Pods
        matching its selector exist. Too few? Create. Too many? Delete.
  \item A \kterm{Deployment} manages ReplicaSets to provide \emph{updates}:
        each new version of the Pod template gets a new ReplicaSet, and the
        Deployment shifts replicas from old to new.
\end{itemize}

\begin{keyidea}
For stateless applications, the \textbf{Deployment} is the object you create.
You will rarely write a ReplicaSet yourself, and almost never a bare Pod.
\end{keyidea}

\section{A complete Deployment}

\begin{yamlcode}{web-deployment.yaml}
apiVersion: apps/v1
kind: Deployment
metadata:
  name: web
  labels:
    app: web
spec:
  replicas: 3
  revisionHistoryLimit: 5          # keep 5 old ReplicaSets for rollback
  selector:
    matchLabels:
      app: web                     # must match template labels
  strategy:
    type: RollingUpdate
    rollingUpdate:
      maxSurge: 1                  # at most 1 extra Pod during the update
      maxUnavailable: 0            # never go below 3 ready Pods
  template:                        # the Pod template
    metadata:
      labels:
        app: web
    spec:
      containers:
        - name: web
          image: nginx:1.27
          ports:
            - containerPort: 80
          readinessProbe:
            httpGet: { path: /, port: 80 }
            periodSeconds: 5
          resources:
            requests: { cpu: 50m, memory: 64Mi }
            limits:   { memory: 128Mi }
\end{yamlcode}

\begin{bashcode}{terminal}
kubectl apply -f web-deployment.yaml
kubectl get deploy,rs,pods -l app=web
kubectl scale deployment web --replicas=5
\end{bashcode}

\begin{tryit}
Delete one of the Pods with \kc{kubectl delete pod <name>} and immediately run
\kc{kubectl get pods -w}. Within a second the ReplicaSet creates a replacement
with a new random suffix. That is self-healing in action.
\end{tryit}

\section{Rolling updates}

Change anything in the Pod template --- the image, an env var, a label --- and
the Deployment performs a \kterm{rolling update}:

\begin{center}
\resizebox{\linewidth}{!}{%
\begin{tikzpicture}[
  v1/.style={regular polygon, regular polygon sides=7, fill=k8sgray!70, text=white,
             minimum size=0.95cm, inner sep=0pt, font=\scriptsize\bfseries},
  v2/.style={regular polygon, regular polygon sides=7, fill=k8sblue, text=white,
             minimum size=0.95cm, inner sep=0pt, font=\scriptsize\bfseries},
  nw/.style={regular polygon, regular polygon sides=7, draw=k8sblue, dashed, text=k8sblue,
             minimum size=0.95cm, inner sep=0pt, font=\scriptsize\bfseries}]
  \foreach \row/\lbl in {0/Start: 3 old Pods,
                         -1.3/Surge: start 1 new Pod,
                         -2.6/{New Pod Ready: remove 1 old, start next},
                         -3.9/Repeat\ldots,
                         -5.2/Done: 3 new Pods}
    \node[font=\footnotesize\bfseries, text=k8snavy, anchor=west] at (-6.2,\row) {\lbl};
  % row 0
  \foreach \x in {0,1.2,2.4} \node[v1] at (\x,0) {v1};
  % row 1
  \foreach \x in {0,1.2,2.4} \node[v1] at (\x,-1.3) {v1};
  \node[nw] at (3.6,-1.3) {v2};
  % row 2
  \foreach \x in {0,1.2} \node[v1] at (\x,-2.6) {v1};
  \node[v2] at (2.4,-2.6) {v2}; \node[nw] at (3.6,-2.6) {v2};
  % row 3
  \node[v1] at (0,-3.9) {v1};
  \foreach \x in {1.2,2.4} \node[v2] at (\x,-3.9) {v2};
  \node[nw] at (3.6,-3.9) {v2};
  % row 4
  \foreach \x in {0,1.2,2.4} \node[v2] at (\x,-5.2) {v2};
  % legend
  \node[v1] at (6.2,-1.3) {v1}; \node[anchor=west, font=\scriptsize] at (6.7,-1.3) {old version};
  \node[v2] at (6.2,-2.6) {v2}; \node[anchor=west, font=\scriptsize] at (6.7,-2.6) {new version, Ready};
  \node[nw] at (6.2,-3.9) {v2}; \node[anchor=west, font=\scriptsize] at (6.7,-3.9) {new version, starting};
\end{tikzpicture}}
\end{center}

With \kc{maxSurge: 1} and \kc{maxUnavailable: 0}, Kubernetes adds one new
Pod, waits until it is \emph{Ready} (readiness probe passes), removes one old
Pod, and repeats. Capacity never drops, so users see no downtime.

\begin{center}
\ktablestart
\begin{tabularx}{\linewidth}{P{3.3cm} L}
\kheader \kth{Setting} & \kth{Effect} \\
\kc{maxSurge} & How many Pods \emph{above} \kc{replicas} may exist during the
                update (number or \%, default 25\%). Higher = faster, more resources. \\
\kc{maxUnavailable} & How many Pods \emph{below} \kc{replicas} may be
                unavailable (default 25\%). 0 = safest. \\
\kc{minReadySeconds} & A new Pod must stay Ready this long before it counts.
                Catches apps that crash shortly after starting. \\
\kc{progressDeadlineSeconds} & If no progress for this long (default 600),
                the rollout is marked failed. \\
\end{tabularx}
\end{center}

\section{Managing rollouts}

\begin{bashcode}{terminal}
# Trigger an update (or edit the YAML and apply it)
kubectl set image deployment/web web=nginx:1.28
kubectl annotate deployment/web kubernetes.io/change-cause="nginx 1.28"

kubectl rollout status deployment/web       # wait and watch progress
kubectl rollout history deployment/web      # list revisions
kubectl rollout history deployment/web --revision=2

kubectl rollout undo deployment/web                 # back to previous revision
kubectl rollout undo deployment/web --to-revision=1

kubectl rollout pause deployment/web        # batch several changes...
kubectl rollout resume deployment/web       # ...then roll them out at once
kubectl rollout restart deployment/web      # re-create Pods (e.g. to reload config)
\end{bashcode}

\begin{tryit}
Set the image to a tag that does not exist: \kc{kubectl set image
deployment/web web=nginx:does-not-exist}. Run \kc{kubectl get pods}: the new
Pod is stuck in \kc{ImagePullBackOff}, but because \kc{maxUnavailable} is 0
your three old Pods keep serving. Fix it with \kc{kubectl rollout undo
deployment/web}.
\end{tryit}

\section{Release strategies}

\begin{center}
\ktablestart
\begin{tabularx}{\linewidth}{P{2.6cm} L P{3.6cm}}
\kheader \kth{Strategy} & \kth{How} & \kth{Trade-off} \\
\textbf{Recreate} & \kc{strategy.type: Recreate}: kill all old Pods, then start new. &
   Downtime; needed when two versions cannot run together. \\
\textbf{Rolling} & Default. Replace Pods gradually. &
   Two versions serve at once for a while. \\
\textbf{Blue/Green} & Run two Deployments (\kc{blue}, \kc{green}); switch the
   Service selector from one to the other. & Instant switch and rollback;
   double resources. \\
\textbf{Canary} & Send a small slice of traffic (say 5\%) to the new version,
   watch metrics, then increase. & Safest for risky changes; needs traffic
   splitting (Gateway API, service mesh, or Argo Rollouts). \\
\end{tabularx}
\end{center}

\begin{protip}
For automated canaries with metric-based promotion and rollback, look at
\textbf{Argo Rollouts} or \textbf{Flagger}. Both build on the concepts in this
chapter and integrate with the Gateway API.
\end{protip}

\begin{recap}
\begin{itemize}
  \item Deployment $\rightarrow$ ReplicaSet $\rightarrow$ Pods.
  \item Changing the Pod template triggers a rolling update via a new ReplicaSet.
  \item Readiness probes + \kc{maxUnavailable: 0} = zero-downtime releases.
  \item \kc{kubectl rollout status/history/undo} manage the release.
\end{itemize}
\end{recap}
